\section{Implementation}
\label{sec:implementation}

We implement \sys around a RISC-V SoC FPGA prototype running Linux and a host
running a modified QEMU device manager. The FPGA and host communicate through a
Xilinx XDMA endpoint. The current logical hierarchy contains two common bridge
functions and up to thirteen downstream-port/endpoint pairs. Software stores 28
configuration-function slots of 4 KiB each in a compact ECAM shadow; a 128-KiB
dual-port BRAM aperture accommodates the image. The DUT accesses a 16-MiB PCI
memory window in which Linux allocates endpoint BARs.

\subsection{FPGA Front-End}

The front-end implements ECAM decoding, direct shadow reads, serialized
configuration writes, and BAR matching. Configuration and BAR events use
32-byte records carried to a host-allocated cyclic DMA ring. A mailbox carries
acknowledgments, BAR read data, route updates, readiness state, and the aggregate
INTx level in the reverse direction. The XDMA bypass interface exposes both raw
DUT memory and a coherent alias used by physical devices and QEMU when reading
or patching queues.

The virtual hierarchy intentionally stops above the PCIe link. Linux uses the
generic ECAM host bridge, while \sys implements neither a physical downstream
link nor PHY, LTSSM, DLLP, or complete TLP behavior.

\subsection{Software Back-End}

The manager owns the XDMA descriptors, compact ECAM image, route table, and
shared interrupt level. Each endpoint object stores its type, virtual BDF,
backend binding, operations table, and private protocol state. Dispatch by BDF
prevents queue, doorbell, and completion state from leaking between endpoints.
Startup fails atomically if any enabled physical backend cannot be acquired,
preventing a logical endpoint from advertising a device that software cannot
operate.

\subsection{NVMe Adapter}

The NVMe adapter maps a physical controller BAR and virtualizes CAP, VS, CC,
CSTS, queue configuration, and doorbells. It maintains independent queue and
command state per endpoint. Before forwarding a submission doorbell, it reads
the DUT SQE through the coherent alias until two consecutive copies agree,
checks basic field validity, and translates metadata, PRP1, PRP2, and any PRP
list entries. Physical CQEs arrive through the same coherent path. The adapter
validates phase, SQ identifier, head, and outstanding command state before
advancing the guest-visible completion and aggregate interrupt.

\subsection{Ethernet Adapter}

The Ethernet adapter binds one Intel 82580 function to one logical endpoint. It
exposes a 512-KiB BAR0 with device ID \texttt{8086:150e} and class
\texttt{020000}. The virtual capability list hides MSI/MSI-X and BAR3, so the
unmodified Linux \texttt{igb} driver selects one legacy TX/RX queue pair. The
adapter shadows queue-0 ring bases and lengths, rewrites descriptor buffer
addresses, forwards TDT/RDT only after descriptor translation, polls physical
writebacks and causes, and emulates virtual ICR/IMS semantics. Controlled reset
stops DMA, masks the physical device, waits for reset, and clears only that
endpoint's virtual state.

\subsection{Current Boundaries}

The implementation supports BAR0-only endpoints and aggregate INTx. It does not
yet support MSI/MSI-X, AER, hotplug, arbitrary switch capabilities, or multiple
NIC queues. Backend configuration occurs before DUT enumeration. These limits
define the claims and experiments in this paper.
